\documentclass[11pt]{article}
\usepackage{nopageno}
\setcounter{secnumdepth}{1}
\usepackage[margin=1in]{geometry}
\begin{document}
\pagestyle{empty}

\subsection{Overview}
\vspace{-0.3cm}

Flood-vulnerable heritage districts dense with unreinforced masonry (URM) construction face the preservation-urgency paradox, the tension between rapid disaster imperatives and technically-informed heritage preservation that causes owners
to diverge from federally recommended infrastructure states with unquantified
consequences for district-scale flood vulnerability. This project couples
social science, hydrodynamic modeling, and spatial data science to investigate
how this paradox operates at the individual building and district scales in two
flood-affected URM heritage districts (Montpelier, VT and Marshall, NC),
test how individual building hardening, district-scale physical interventions,
and mixed approaches compare in their effects on district-scale flood vulnerability,
and determine what organizational and network structures produce the coordination
failures that prevent communities from pursuing district-scale intervention. 

\vspace{-0.5cm}

\subsection{Intellectual Merit}
\vspace{-0.3cm}

This research produces the first empirical comparison of individual building
hardening, district-scale physical interventions, and mixed approaches in
their collective flood vulnerability outcomes, and the first characterization
of the organizational and network structures that explain why communities
default to individual hardening rather than coordinated district-scale
intervention, advancing understanding of how collective action structure
shapes infrastructure system resilience and extending NIST's socio-technical
systems framing to flood hazards and heritage infrastructure. RQ1 establishes
how the preservation-urgency paradox operates at individual and district
scales and which buildings and owners are drawn into collective action around
building-level hardening or district-scale intervention versus acting independently, using mixed-methods data collection, social network
analysis, and feature importance modeling across two contrasting sites. RQ2 tests how individual building hardening,
district-scale physical interventions, and mixed approaches compare in their
effects on district-scale flood vulnerability using HEC-RAS 2D simulations, and whether any intervention produces threshold
behavior in collective flood outcomes using PELT change-point detection and Sobol index
sensitivity analysis via Gaussian process emulation. RQ3 characterizes the
organizational and network structures that explain exclusion from collective action
and whether those same structures explain why communities defaulted to fragmented
individual building hardening rather than district-scale intervention where technical
expertise identified it as an option, giving communities a quantitative basis for
pre-disaster investment decisions.

\vspace{-0.5cm}

\subsection{Broader Impacts}

\vspace{-0.3cm}

The project delivers the first empirical comparison of building-level and
district-scale intervention approaches in URM heritage districts and the
first characterization of the organizational dynamics that explain coordination
failure around district-scale intervention, giving communities a quantitative
basis for pre-disaster investment that currently has none. The
Flood-Resilient Heritage Preservation Toolkit, developed through the Association for Preservation Technology (APT) Disaster Response Initiative (DRI) and reviewed by Jennifer Wellock (Department of Interior (DOI) National Coordinator Natural and Cultural Resources Recovery Support Function), serves preservation professionals,
property owners, and community leaders. Distribution reaches heritage-at-risk communities nationwide through the APT network, federal
heritage recovery channels. The methodology is
transferable to other hazards where individual decisions aggregate to
collective infrastructure outcomes, contributing to both NEHRP and NWIRP
goals. All data and code are archived openly in the NHERI DesignSafe Data
Depot. Research outcomes integrate into education through PI Napolitano's and Co-PI Busse's Penn State and Co-PI Chi's Indiana University courses, with postdoctoral
scholars trained in spatial data science, hydrodynamic modeling, and social science. Undergraduate research assistants will
gain experience in building typology classification, geospatial data preparation, and hydrodynamic output processing through PI Napolitano's role as undergraduate research
coordinator for her department.

\end{document}

